% !TeX program = xelatex
% ============================================================================
%  第 13 课　它有什么用　—— 学生版讲义
%  与 02-课件/第13课-它有什么用/第13课-它有什么用.tex 同步
% ============================================================================

\xhlesson{13}{它有什么用}{}

\begin{mubiao}
  \begin{itemize}
    \item 会用费马小定理快速求大幂的余数。
    \item 能回答第 1 课留下的问题：\(999999\) 为什么能被 \(7\) 整除。
    \item 知道身份证校验码怎么算，RSA 大致在做什么。
  \end{itemize}
\end{mubiao}

\section*{一、回答第 1 课的问题}

\noindent 第 1 课，我们在黑板最右边写下 \(999999\div7=?\)。今天，把它划掉。

\begin{dongshou}[一步一步写]
  第一步，改写：
  \[
    999999=1000000-1=10^{6}-\underline{\hspace{2em}} .
  \]
  第二步，用定理。\(7\) 是素数，\(10\) 不是 \(7\) 的倍数，取 \(p=7\)、\(a=10\)：
  \[
    10^{\,7-1}=10^{6}\equiv\underline{\hspace{2em}} \pmod 7 .
  \]
  第三步，代回去：
  \[
    10^{6}-1\equiv\underline{\hspace{2em}}-\underline{\hspace{2em}}
    =\underline{\hspace{2em}} \pmod 7 .
  \]
\end{dongshou}

\begin{sikao}
  余数是 \(0\)，说明什么？\\
  从头到尾，我们做过一次除法、用过一次计算器吗？
  \liukong{2}
\end{sikao}

\section*{二、快速求大幂的余数}

\begin{example}[\(2^{100}\) 除以 \(7\)]
  不用把 \(2^{100}\) 算出来，看指数：
  \[
    100=6\times16+4 .
  \]
  由费马小定理 \(2^{6}\equiv1\pmod 7\)，于是
  \[
    2^{100}=2^{6\times16+4}=(2^{6})^{16}\times2^{4}
    \equiv1^{16}\times16\equiv16\equiv2 \pmod 7 .
  \]
\end{example}

\begin{dongshou}[自己练：\(3^{200}\) 除以 \(11\)]
  \(11\) 是素数，由费马小定理 \(3^{10}\equiv1\pmod{11}\)。
  \[
    200=10\times\underline{\hspace{1.6em}},
    \qquad
    3^{200}=(3^{10})^{\underline{\hspace{1.6em}}}
    \equiv\underline{\hspace{1.6em}}^{\,\underline{\hspace{1.2em}}}
    =\underline{\hspace{1.6em}} \pmod{11} .
  \]
\end{dongshou}

\begin{sikao}
  求 \(5^{2026}\) 除以 \(7\) 的余数，第一步应该看什么？\\
  提示：\(5^{6}\equiv1\pmod 7\)，先算 \(2026\div6\) 的余数。
  \liukong{3}
\end{sikao}

\section*{三、身份证最后一位，是怎么算出来的}

\noindent 身份证前 \(17\) 位，每一位乘一个固定的\textbf{权重}，加起来，再除以 \(11\) 取余；
余数查表换成最后一位。

\begin{center}
\begin{tabular}{@{}lc*{17}{c}@{}}
  \toprule
  位数 & \(1\) & \(2\) & \(3\) & \(4\) & \(5\) & \(6\) & \(7\) & \(8\) & \(9\) & \(10\) & \(11\) & \(12\) & \(13\) & \(14\) & \(15\) & \(16\) & \(17\) \\
  \midrule
  权重 & \(7\) & \(9\) & \(10\) & \(5\) & \(8\) & \(4\) & \(2\) & \(1\) & \(6\) & \(3\) & \(7\) & \(9\) & \(10\) & \(5\) & \(8\) & \(4\) & \(2\) \\
  \bottomrule
\end{tabular}
\end{center}

\begin{center}
\begin{tabular}{@{}l*{11}{c}@{}}
  \toprule
  余数 & \(0\) & \(1\) & \(2\) & \(3\) & \(4\) & \(5\) & \(6\) & \(7\) & \(8\) & \(9\) & \(10\) \\
  \midrule
  校验码 & \(1\) & \(0\) & X & \(9\) & \(8\) & \(7\) & \(6\) & \(5\) & \(4\) & \(3\) & \(2\) \\
  \bottomrule
\end{tabular}
\end{center}

\begin{dongshou}[跟着老师算一遍]
  取前 \(17\) 位 \(3\,5\,0\,6\,0\,4\,2\,0\,1\,3\,0\,5\,2\,0\,1\,2\,3\)（虚构示例），
  每一位乘上对应权重再相加：
  \[
    21+45+0+30+0+16+4+0+6+9+0+45+20+0+8+8+6=\underline{\hspace{3em}} .
  \]
  \[
    218=11\times\underline{\hspace{1.4em}}+\underline{\hspace{1.4em}}
    \quad\Longrightarrow\quad \text{余数是 }\underline{\hspace{1.4em}}
    \quad\Longrightarrow\quad \text{校验码是 }\underline{\hspace{1.4em}} .
  \]
\end{dongshou}

\begin{sikao}
  为什么改掉前 \(17\) 位里的任意一位，机器就能发现号码不对？
  \liukong{3}
\end{sikao}

\section*{四、RSA：现在最有名的用处}

\noindent 上网的时候，浏览器和网站之间要先悄悄商量一把「钥匙」，靠的就是 RSA。
它做的是这样一件事：

\begin{itemize}
  \item 有一把\textbf{公钥}，像一把谁都能拿到的「锁」：谁都可以用它把消息锁上；
  \item 有一把\textbf{私钥}，像一把只有收信人才有的钥匙：只有它能打开锁上的消息。
\end{itemize}

\noindent 别人半路偷看也没用：公钥只能锁、不能开。
而钥匙是用两个很大的素数\emph{乘}出来的——想反过来把一个大数\emph{拆回}两个素数，非常难。

\begin{sikao}
  课上的例子是 \(3\times11=33\)，小到一眼就能拆开。\\
  如果换成两个五十位的大素数相乘呢？想一想：拆开它，难在哪里？
  \liukong{3}
\end{sikao}

\noindent\textbf{一点数学史}：\(1977\) 年，美国的 Rivest、Shamir、Adleman 三人提出这套方法，
名字就取三人姓氏的首字母，叫 RSA；\(2002\) 年，他们因此获得了图灵奖。

\begin{xiaojie}
  这节课要带走的三句话：
  \begin{itemize}
    \item 费马小定理回答了 \(999999\div7\)：\(10^{6}\equiv1\pmod 7\)，所以 \(7\) 整除它。
    \item 它还能把巨大的幂「缩」小：\(2^{100}\equiv2\pmod 7\)。
    \item 身份证校验码、手机里的加密，用的都是这套「取余」的算术。
  \end{itemize}
\end{xiaojie}

\noindent\textbf{下节课}：复习、总结与展望——把整门课串成一条线。